\section{Build flow and traceability}
\label{sec:flow}

\subsection{One JSON file per build}
Every build starts from a JSON file in \file{configs/gemm/}. \Cref{lst:config} is the best configuration. The
\code{kernel} section holds every \code{gemm\_hls} parameter; \code{memory\_map} wires the matrices to channels;
\code{clocks} sets the accel clock and the HLS clock uncertainty; \code{vivado} picks a preset and lets any directive
be overridden; \code{max\_threads: auto} gives Vivado \(\min(8, \text{cores})\) threads.

\lstinputlisting[language={}, caption={The best configuration, \file{configs/gemm/pe32x8\_ta\_t1024x512\_225mhz.json}.},
  label={lst:config}, basicstyle=\ttfamily\scriptsize]{../../configs/gemm/pe32x8_ta_t1024x512_225mhz.json}

\code{BuildCfg} (\file{flow/gemmflow/config.py}) loads the file into dataclasses, rejects unknown keys, and checks
every rule that \code{gemm\_hls} enforces in CMake or with \code{static\_assert}s before any tool runs: bus widths
against the element size, \(P_M\) against the M bus word, tiles against \(P_N\), \(P_M\) and the bus words, the
double-buffering rule \(T_M/P_M \ge T_N/P_N\), the transposed-A rule and more. A violation is a configuration error
with a readable message, not a failed HLS run twenty minutes later. It also renders the kernel's \code{Config.h} and
computes three hashes:
\begin{itemize}
  \item \code{config\_hash}: SHA-256 of the whole resolved configuration. The build id is
        \code{<name>-<config\_hash[:8]>}, and the first 32 bits are compiled into the bitstream.
  \item \code{hls\_hash}: the kernel, clock and HLS sections plus a hash of the kernel source tree. Builds that differ
        only in Vivado settings share one HLS result.
  \item the stamps' input hashes (below).
\end{itemize}
Names follow the parameters (\code{f32\_pe32x8\_t1024x512\_225mhz\_tA\_aggressive}), so the same point in two
campaigns, or a campaign point and a hand-written config with the same name, is the same build directory.

\subsection{Steps, stamps and outputs}
\Cref{fig:flow} shows the steps. Each one hashes its inputs and writes a stamp next to its outputs; if the stamp
matches and the outputs exist, the step is skipped:
\begin{description}[font=\sffamily\small\bfseries, leftmargin=1.2em]
  \item[\code{HlsStep}] runs \code{vitis-run --mode hls} with \file{tcl/hls/run\_hls.tcl} (C synthesis, IP export).
        \code{vitis-run} exits 0 even when synthesis fails, so the flow scans the log for errors. Several HLS runs
        can go in parallel.
  \item[\code{SocStep}] runs \file{hw/gemm/gemm\_soc.py}: LiteX generates the Verilog and runs its own Vivado build
        script with our hooks (IP import before synthesis) and the resolved directives. Its inputs include the source
        trees of \file{hw/gemm}, \file{boards/dupcb} and \file{tcl/vivado}.
  \item[\code{HostStep}] builds \code{liblitepcie} and the host tools against the build's \code{csr.h}.
\end{description}
A build directory holds the resolved config, \code{manifest.json} (git commit and dirty flag, tool versions, host),
\code{reports/summary.json} (timing per check, utilisation of the whole design and of the kernel, HLS estimates) and
the logs. With 15\,GB of RAM on the build machine, one Vivado run of the 480T at a time is the limit; a PE 32\x8 build
takes 25--35 minutes, a PE 8\x8 build 12--14.

\begin{figure}[t]
  \centering
  \resizebox{\linewidth}{!}{% Build flow from one JSON config to results, with the cached steps and their outputs.
\begin{tikzpicture}[x=1cm,y=1cm, every node/.append style={font=\sffamily\footnotesize}]
  \tikzset{st/.style={flw, minimum width=30mm, minimum height=11mm, align=center},
           outp/.style={art, text width=30mm},
           stamp/.style={draw=cGreen!60!black, fill=cGreen!10, font=\sffamily\tiny, rounded corners=1pt, inner sep=1.5pt}}
  % row 1: config -> HLS -> SoC -> host
  \node[art, text width=27mm, minimum height=11mm] (cfg) at (0,0)
       {\code{configs/gemm/*.json}\\ kernel, memory map, clocks, HLS, Vivado, SoC};
  \node[st] (bc) at (3.6,0) {\code{BuildCfg}\\[-1pt]\scriptsize validate rules, hashes,\\[-1pt]\scriptsize render \code{Config.h}};
  \node[st] (hls) at (7.4,0) {\code{HlsStep}\\[-1pt]\scriptsize \code{vitis-run} + \file{run\_hls.tcl}\\[-1pt]\scriptsize csynth, IP export};
  \node[st] (soc) at (11.2,0) {\code{SocStep}\\[-1pt]\scriptsize \file{gemm\_soc.py} (LiteX)\\[-1pt]\scriptsize + Vivado, Tcl hooks};
  \node[st] (hst) at (15.0,0) {\code{HostStep}\\[-1pt]\scriptsize liblitepcie,\\[-1pt]\scriptsize \code{gemm\_runner} vs \code{csr.h}};
  \draw[link] (cfg) -- (bc); \draw[link] (bc) -- (hls); \draw[link] (hls) -- (soc); \draw[link] (soc) -- (hst);
  % stamps
  \foreach \n in {hls,soc,hst} { \node[stamp, anchor=south east] at (\n.north east) {stamp: inputs hash}; }
  % outputs
  \node[outp] (o1) at (7.4,-1.9) {\file{build/gemm/hls/<hls\_hash>/}\\ IP catalog, \code{csynth.xml}; shared by all Vivado knobs};
  \node[outp] (o2) at (11.2,-1.9) {bitstream, \code{csr.csv}/\code{csr.h}, LitePCIe driver, timing + utilisation reports};
  \node[outp] (o3) at (15.0,-1.9) {\file{<build>/sw/}: \code{gemm\_runner}, \code{ddr\_dma\_test}};
  \foreach \a/\b in {hls/o1, soc/o2, hst/o3} { \draw[->, black!50] (\a) -- (\b); }
  \node[art, text width=60mm, minimum height=9mm] (bd) at (11.2,-3.55)
       {\file{build/gemm/soc/<name>-<config\_hash8>/}: \code{config.json}, \code{manifest.json} (git, tools),
        \code{reports/summary.json}, \code{stamps/}, \code{logs/}};
  \draw[->, black!50] (o2) -- (bd); \draw[->, black!50] (o3.south) |- (bd.east);
  % Vivado detail
  \node[blk, fill=black!4, align=left, text width=44mm, font=\sffamily\scriptsize, anchor=north west] (viv) at (-1.35,-1.3)
       {\textbf{Vivado in \code{SocStep}} (LiteX build script)\\
        \code{pre\_synthesis\_commands}: \file{import\_hls\_ip.tcl} (\code{ip\_repo\_paths}, \code{create\_ip},
        \code{generate\_target}, \code{synth\_ip})\\
        directives from the preset (default / explore / aggressive), \code{maxThreads}\\
        The kernel \code{Instance} is generated from the HLS port list.};
  % row 3: board
  \node[st, fill=dHost] (prog) at (0.6,-5.6)  {program\\[-1pt]\scriptsize JTAG, \file{program.tcl}};
  \node[st, fill=dHost] (scan) at (4.3,-5.6)  {PCIe rescan};
  \node[st, fill=dHost] (drv)  at (8.0,-5.6)  {driver build + load\\[-1pt]\scriptsize probe resets the SoC};
  \node[st, fill=dHost] (ddr)  at (11.7,-5.6) {DDR calibration\\[-1pt]\scriptsize \file{ddr\_init.sh}, both channels};
  \node[st, fill=dHost] (run)  at (15.4,-5.6) {\code{gemm\_runner}\\[-1pt]\scriptsize sizes \x repetitions};
  \draw[link] (prog) -- (scan); \draw[link] (scan) -- (drv); \draw[link] (drv) -- (ddr); \draw[link] (ddr) -- (run);
  \draw[link] (bd.south) -- ++(0,-0.45) -| (prog.north);
  \node[art, text width=33mm, minimum height=9mm] (res) at (15.4,-7.3)
       {\file{results/}: JSONL (source of truth) + CSV; ids, hashes, git, timing, counters};
  \draw[link] (run) -- (res);
  \node[lbl, anchor=south west] at ([xshift=2mm]prog.north) {\code{Board.bring\_up()} (\file{flow/gemmflow/board.py}), in this order};
  \node[note, anchor=north west, text width=88mm, align=left] at (-1.35,-6.55)
       {\code{scripts/gemm\_build.py <config>} runs the top row (unchanged steps are skipped);
        \code{scripts/gemm\_run.py <build> --program} runs the bottom row;
        \code{scripts/gemm\_dse.py <campaign>} does both for every point of a design space
        (\cref{sec:dse}). The DDR calibration must come after the driver load, because the driver's probe resets the SoC.};
\end{tikzpicture}
}
  \caption{From a JSON configuration to results. Top: the cached build steps and what they produce. Bottom: the
    bring-up on the card, in the only order that works, and the runs.}
  \label{fig:flow}
\end{figure}

\subsection{Vivado strategy presets}
\Cref{tab:presets} lists the presets. Explicit directives in the config override the preset. Vivado's \emph{default}
strategy was enough for most builds; \cref{sec:eval:timing} shows where the others mattered.

\begin{table}[ht]
\centering\scriptsize\setlength{\tabcolsep}{4pt}
\caption{Vivado presets (\code{VIVADO\_PRESETS} in \file{flow/gemmflow/config.py}).}
\label{tab:presets}
\begin{tabular}{@{}lllllll@{}}
\toprule
Preset & synth & opt & place & phys\_opt (place) & route & phys\_opt (route) \\
\midrule
default & default & default & default & -- & default & default \\
explore & default & Explore & Explore & Explore & Explore & Explore \\
aggressive & PerformanceOptimized & ExploreWithRemap & ExtraTimingOpt & AggressiveExplore & AggressiveExplore & AggressiveExplore \\
\bottomrule
\end{tabular}
\end{table}

\subsection{Tests}
The Python flow has unit tests (\file{flow/tests/}, standard \code{unittest}, under a second, no tools or card): every
configuration rule with its message, size rules, names, presets, \code{Config.h} rendering, the result store, campaign
expansion, build reports, and the build ids of every committed configuration. The last one matters because build reuse
depends on stable hashes: a new config field or a changed default would otherwise silently force full rebuilds. The
DMA bridge and the C write converter have Migen simulations (\file{hw/gemm/tests/}).
